\documentclass[11pt]{article}

\usepackage[margin=1in]{geometry}
\usepackage{amsmath,amssymb,amsthm,mathtools}
\usepackage[T1]{fontenc}
\usepackage[utf8]{inputenc}
\usepackage{lmodern}
\usepackage{microtype}
\usepackage{enumitem}
\usepackage{booktabs,array,longtable}
\usepackage{verbatim}
\usepackage[colorlinks=true,linkcolor=blue,citecolor=blue,urlcolor=blue]{hyperref}

\allowdisplaybreaks[2]
\setlist[itemize]{leftmargin=1.5em,itemsep=0.25em,topsep=0.35em}
\setlist[enumerate]{leftmargin=1.5em,itemsep=0.25em,topsep=0.35em}
\emergencystretch=2em

\newcommand{\E}{\mathbb{E}}
\newcommand{\R}{\mathbb{R}}
\newcommand{\cF}{\mathcal{F}}
\newcommand{\cP}{\mathcal{P}}
\newcommand{\cZ}{\mathcal{Z}}
\newcommand{\cE}{\mathcal{E}}
\newcommand{\T}{\top}
\newcommand{\PC}{\mathrm{PC}}
\newcommand{\SDF}{\mathrm{SDF}}
\newcommand{\Cov}{\operatorname{Cov}}
\newcommand{\Var}{\operatorname{Var}}
\newcommand{\Corr}{\operatorname{Corr}}
\newcommand{\rank}{\operatorname{rank}}
\newcommand{\diag}{\operatorname{diag}}
\newcommand{\sgn}{\operatorname{sgn}}
\newcommand{\plim}{\operatorname*{plim}}
\newcommand{\kurt}{\operatorname{kurt}}
\newcommand{\se}{\operatorname{se}}

\theoremstyle{plain}
\newtheorem{proposition}{Proposition}
\newtheorem{theorem}{Theorem}
\theoremstyle{definition}
\newtheorem{assumption}{Assumption}
\theoremstyle{remark}
\newtheorem{remark}{Remark}

\newcommand{\heading}[1]{\vspace{0.5em}\noindent\textbf{#1}}

\title{Identification, Inference, and R Bootstrap Options for a Scalar Triangular System\\
\large A consolidated self-contained guide for point-identified and set-identified cases}
\author{}
\date{June 2026}

\begin{document}
\maketitle


\begin{abstract}
This consolidated report combines the earlier identification guide for a scalar
triangular system with a practical inference and implementation guide for
standard errors, confidence intervals, and bootstrap methods in R.  The maintained
empirical object is a linear triangular model with one outcome, one endogenous
regressor, one scalar external or constructed instrument when such a variable is
used, and a common set of controls.  The first part translates the identification
ideas in Lewbel (2012), Lee--Lewbel--Schennach--Zhang's \texttt{trigmm} paper,
and Lewis's survey on higher-moment identification into this scalar setting.  The
second part focuses on the inferential problem: how to bootstrap the
point-identified generated-instrument estimator, how to bootstrap scalar
set-identified endpoints, how to reproduce a \texttt{trigmmset}-style second
moment confidence construction in R, and which R packages provide reusable
building blocks.  The document is written to be usable without prior familiarity
with partial identification, while retaining the formulas needed to implement the
methods directly.
\end{abstract}

\tableofcontents

\section{Executive summary}
\label{sec:exec}

The target coefficient is the scalar structural effect of the endogenous
regressor, denoted by $\theta$. The most useful summary is the following.

\begin{longtable}{>{\raggedright\arraybackslash}p{0.18\textwidth}>
                  {\raggedright\arraybackslash}p{0.25\textwidth}>
                  {\raggedright\arraybackslash}p{0.25\textwidth}>
                  {\raggedright\arraybackslash}p{0.22\textwidth}}
\toprule
Option & Main identifying information & What is identified in the scalar triangular model & Main practical cautions \\
\midrule
\endfirsthead
\toprule
Option & Main identifying information & What is identified in the scalar triangular model & Main practical cautions \\
\midrule
\endhead
Lewbel exact heteroskedasticity & A scalar $Z$ satisfies
$\Cov(Z,\varepsilon_1\varepsilon_2)=0$ and
$\Cov(Z,\varepsilon_2^2)\ne0$. & Point identification:
$\theta=\Cov(Z,W_1W_2)/\Cov(Z,W_2^2)$. & Needs heteroskedastic relevance of the first-stage residual. With one instrument it is just identified, so overidentification tests are unavailable. \\
\addlinespace
Lewbel bounded-invalidity set & Same scalar $Z$, but
$\Corr(Z,\varepsilon_1\varepsilon_2)$ is allowed to be nonzero, bounded relative to $\Corr(Z,\varepsilon_2^2)$. & A finite interval for $\theta$ when the relevance covariance is nonzero and the slack is below one. & The bound depends on a user-chosen slack $\tau$; sampling uncertainty requires an outer confidence construction or bootstrap. \\
\addlinespace
Lee--Lewbel--Schennach--Zhang point identification (\texttt{trigmm}) & Common-factor errors
$W_2=u+v$, $W_1=\alpha u+\theta v+r$ with independent components and enough non-Gaussianity. & Point identification of $(\alpha,\theta)$ from two higher-cumulant equations $M_p(\alpha,\theta)=0$; then all other coefficients follow from the common residualization map. & Requires the sign of the common-factor loading and nondegenerate higher cumulants. Nonlinear GMM can have local minima and weak higher-moment identification. \\
\addlinespace
Lee--Lewbel--Schennach--Zhang sharp set & Common-factor structure without the non-Gaussianity point-identifying condition; uses the implied decomposition of observables into Gaussian and non-Gaussian independent parts. & A sharp population joint set for the latent common-factor parameters. & Conceptually important but hard to estimate and interpret; the Stata paper therefore reports a simpler covariance-based outer bound. \\
\addlinespace
Lee--Lewbel--Schennach--Zhang set identification (\texttt{trigmmset}) & Same common-factor structure and sign of common-factor loading, but no non-Gaussianity; only second moments. & A joint set in $(\alpha,\theta)$. For positive common-factor loading, $\theta\le B_0\le\alpha$ and $(B_0-\theta)(\alpha-B_0)\le D_0$. The projection onto $\theta$ alone is generally one-sided, not finite. & The reported confidence set is conservative and is a joint region; a finite interval for $\theta$ alone needs an extra magnitude restriction on the common-factor loading. \\
\addlinespace
Lee--Lewbel--Schennach--Zhang plus an external instrument & Add a conventional IV moment such as $\E[Z(W_1-\theta W_2)]=0$ to the internal non-Gaussian GMM moments. & Usually point identification, often overidentification, with a direct comparison between the external-IV and internal-moment estimates. & If the two information sources conflict, the overidentifying restrictions can reject but cannot by themselves say which assumption failed. \\
\addlinespace
Variance-regime / Rigobon-style identification & A scalar state or instrument $Z$ partitions the sample into variance regimes; the structural impact matrix is fixed while shock variances change nonproportionally. & In a square two-shock representation, the impact ratio for the labeled $W_2$ shock can be identified from regime covariance matrices; in the one-changing-shock case it is a simple covariance-difference ratio. & This is an alternative square-shock model, not the three-shock common-factor model. Needs a credible regime/state and shock labeling. \\
\addlinespace
Parametric volatility and unconditional volatility moments & GARCH, Markov-switching, smooth transition, stochastic-volatility, or autocovariances of squared residuals provide additional covariance equations. & Identifies a square impact matrix and hence an impact ratio after normalization and labeling. & Strong functional-form or time-series assumptions; high-moment estimates can be noisy. \\
\addlinespace
Non-Gaussian ICA / likelihood / higher-moment GMM & Independent or partly independent non-Gaussian shocks in a square decomposition of $(W_2,W_1)$. & Identifies the square mixing matrix up to sign, scale, and column order; $\theta$ is an impact ratio for the correctly labeled column. & Needs shock labeling, sign normalization, and enough departure from Gaussianity. Standard Wald intervals are unreliable under weak non-Gaussianity. \\
\bottomrule
\end{longtable}

The bottom line is that there are two direct methods for the exact Lewbel
triangular regression problem. Lewbel's method uses the scalar instrument $Z$ to
construct a heteroskedasticity-based instrument. Lee--Lewbel--Schennach--Zhang
instead do not need $Z$ for their internal point or set identification; they use a
common-factor decomposition of the errors. The methods surveyed by Lewis are
valuable alternatives and diagnostics, but they become direct estimators of
$\theta$ only after one changes the latent structure to a square two-shock
representation and labels one shock as the shock whose effect on $W_1$ per unit
of $W_2$ is the desired coefficient.

\section{Target scalar triangular system and notation}
\label{sec:model}

The notation follows the supplied TeX note, specialized to one endogenous
regressor. Let $Y_{1,t+1}$ be the outcome and let $Y_{2,t+1}$ be a scalar
endogenous regressor. Let $X_t$ be the common vector of controls used in both the
outcome equation and the reduced form for the endogenous regressor. The
triangular model is
\begin{align}
Y_{1,t+1} &= X_t^\T\beta_1 + \theta Y_{2,t+1}+\varepsilon_{1,t+1},
\label{eq:structural}\\
Y_{2,t+1} &= X_t^\T\beta_2^R+\varepsilon_{2,t+1}.
\label{eq:firststage}
\end{align}
The coefficient of interest is the scalar $\theta$. The errors may be correlated,
so ordinary least squares in \eqref{eq:structural} is not generally causal.

Project both observed variables on the same $X_t$ and define the residuals
\begin{align}
W_{1,t+1} &:= Y_{1,t+1}-X_t^\T\beta_1^R,\qquad
\beta_1^R:=\bigl(\E[X_tX_t^\T]\bigr)^{-1}\E[X_tY_{1,t+1}],\label{eq:W1}\\
W_{2,t+1} &:= Y_{2,t+1}-X_t^\T\beta_2^R,
\qquad
\beta_2^R:=\bigl(\E[X_tX_t^\T]\bigr)^{-1}\E[X_tY_{2,t+1}].\label{eq:W2}
\end{align}
Under $\E[X_t\varepsilon_{1,t+1}]=\E[X_t\varepsilon_{2,t+1}]=0$, the residualized
system is
\begin{equation}
W_{1,t+1}=\theta W_{2,t+1}+\varepsilon_{1,t+1},
\qquad
W_{2,t+1}=\varepsilon_{2,t+1}.
\label{eq:resid-system}
\end{equation}
For any candidate value $w$ of $\theta$, the implied structural residual is
\begin{equation}
\varepsilon_{1,t+1}(w):=W_{1,t+1}-wW_{2,t+1}.
\label{eq:cand-resid}
\end{equation}
Once $\theta$ is known or restricted to a set $\Theta$, the non-endogenous
coefficients in the outcome equation are recovered by the same linear map used
in the supplied note:
\begin{equation}
\beta_1(w)=\beta_1^R-\beta_2^R w.
\label{eq:beta-recovery}
\end{equation}
Thus each candidate $w$ for $\theta$ implies one and only one candidate
coefficient vector $\beta_1(w)$. This is why a finite interval for $\theta$ is the
central object in the scalar problem.

Throughout, $Z_{t+1}$ denotes a single observed scalar that may be used as an
external instrument, a heteroskedasticity-generating variable, a variance-regime
indicator, or a shock-labeling variable depending on the identification strategy.
The role of $Z_{t+1}$ is strategy-specific; it is not automatically a valid
instrument simply because it is observed.

\section{Baseline: Lewbel's scalar heteroskedasticity method}
\label{sec:lewbelbaseline}

Although the two new attached papers are the main focus, the Lewbel baseline is
useful because it fixes the notation and clarifies what it means to have one
endogenous regressor and one instrument.

\subsection{Exact validity gives a point estimate}

Lewbel's exact triangular condition is
\begin{equation}
\Cov(Z_{t+1},\varepsilon_{1,t+1}\varepsilon_{2,t+1})=0,
\qquad
\Cov(Z_{t+1},\varepsilon_{2,t+1}^2)\ne0.
\label{eq:lewbel-exact-assump}
\end{equation}
The first condition says that $Z_{t+1}$ is not correlated with the product of
the two structural errors. The second condition says that $Z_{t+1}$ predicts the
variance of the first-stage error. Therefore
\begin{equation}
H_{t+1}:=(Z_{t+1}-\E Z_{t+1})W_{2,t+1}
\label{eq:lewbel-instrument}
\end{equation}
is a valid constructed instrument for $W_{2,t+1}$ in the residualized equation.
Indeed,
\begin{align}
\E[H_{t+1}\varepsilon_{1,t+1}]&=
\Cov(Z_{t+1},\varepsilon_{1,t+1}\varepsilon_{2,t+1})=0,\label{eq:lewbel-valid}\\
\E[H_{t+1}W_{2,t+1}]&=\Cov(Z_{t+1},\varepsilon_{2,t+1}^2)\ne0.\label{eq:lewbel-relevance}
\end{align}
In the scalar residualized system this yields the closed-form population
coefficient
\begin{equation}
\theta=
\frac{\Cov(Z_{t+1},W_{1,t+1}W_{2,t+1})}
     {\Cov(Z_{t+1},W_{2,t+1}^{2})}.
\label{eq:lewbel-point}
\end{equation}
With controls, the same calculation is applied after residualizing both
$Y_{1,t+1}$ and $Y_{2,t+1}$ on $X_t$, and the remaining coefficients follow from
\eqref{eq:beta-recovery}.

\subsection{Bounded invalidity gives a scalar interval}

The supplied TeX note generalizes the exact condition by allowing a bounded
amount of contamination. In the scalar case, assume that for a known
$\tau\in[0,1)$,
\begin{equation}
\left|\Corr(Z_{t+1},\varepsilon_{1,t+1}\varepsilon_{2,t+1})\right|
\le
\tau\left|\Corr(Z_{t+1},\varepsilon_{2,t+1}^2)\right|.
\label{eq:relative-corr}
\end{equation}
Define the observable population quantities
\begin{align}
L&:=\Cov(Z_{t+1},W_{1,t+1}W_{2,t+1}),&
Q&:=\Cov(Z_{t+1},W_{2,t+1}^2),\label{eq:LQ}\\
S_0&:=\Var(W_{1,t+1}W_{2,t+1}),&
S_1&:=\Cov(W_{2,t+1}^2,W_{1,t+1}W_{2,t+1}),&
S_2&:=\Var(W_{2,t+1}^2).
\label{eq:S012}
\end{align}
For a candidate $w$, the relative-correlation bound becomes the quadratic
inequality
\begin{equation}
(L-Qw)^2-d\{S_0-2S_1w+S_2w^2\}\le0,
\qquad d:=\tau^2\frac{Q^2}{S_2}.
\label{eq:scalar-lewbel-set}
\end{equation}
If $Q\ne0$ and $\tau<1$, the leading coefficient is
$Q^2-dS_2=Q^2(1-\tau^2)>0$, so the set is a closed finite interval whenever the
discriminant is nonnegative. Writing
\begin{equation}
a:=Q^2-dS_2,\\qquad
b:=-2LQ+2dS_1,
\qquad
c:=L^2-dS_0,
\label{eq:abc-lewbel}
\end{equation}
the interval is
\begin{equation}
\Theta_{\text{Lewbel}}(\tau)=
\left[
\frac{-b-\sqrt{b^2-4ac}}{2a},\,
\frac{-b+\sqrt{b^2-4ac}}{2a}
\right].
\label{eq:lewbel-interval}
\end{equation}
At $\tau=0$, the interval collapses to the point in \eqref{eq:lewbel-point}. As
$\tau$ grows, the interval widens. The interpretation is simple: the more
invalidity one is willing to tolerate in the constructed instrument, the less
precisely $\theta$ is identified.

\subsection{Sampling uncertainty}

The exact Lewbel point estimator can be computed by two-stage least squares using
$X_t$ and $(Z_{t+1}-\bar Z)\widehat W_{2,t+1}$ as instruments. Heteroskedasticity-
robust or cluster-robust covariance matrices are standard. For the interval
\eqref{eq:lewbel-interval}, a conservative confidence set can be obtained by
replacing the population moments with sample moments and using a bootstrap or a
moment-inequality outer construction. The attached Lee et al. paper provides a
more explicit closed-form confidence construction for its own second-moment set,
which is described in Section~\ref{sec:lsz-set}.

\section{Direct option 1 from Lee--Lewbel--Schennach--Zhang: point identification by non-Gaussianity}
\label{sec:lsz-point}

The \texttt{trigmm} paper studies the same triangular shape as
\eqref{eq:structural}--\eqref{eq:firststage}, but it does not use a
Lewbel-style scalar $Z$ to construct an instrument. Instead, it imposes a latent
common-factor structure on the residuals. In the notation of this report,
\begin{equation}
W_{2,t+1}=u_{t+1}+v_{t+1},
\qquad
W_{1,t+1}=\theta W_{2,t+1}+\beta_u u_{t+1}+r_{t+1}.
\label{eq:lsz-raw}
\end{equation}
Equivalently,
\begin{equation}
W_{1,t+1}=\alpha u_{t+1}+\theta v_{t+1}+r_{t+1},
\qquad
\alpha:=\theta+\beta_u.
\label{eq:alpha-def}
\end{equation}
Here $u$ is the unobserved common cause that creates endogeneity, $v$ is the
idiosyncratic first-stage shock, and $r$ is the idiosyncratic outcome shock. The
key assumptions are:
\begin{itemize}
\item $u_{t+1}$, $v_{t+1}$, and $r_{t+1}$ are mutually independent, and in the
implemented command they are also independent of $X_t$;
\item the sign of $\beta_u=\alpha-\theta$ is specified by the researcher;
\item sufficiently high moments exist and the relevant higher cumulants are
nondegenerate.
\end{itemize}
The sign condition matters because replacing $u$ by $-u$ changes the sign of the
common-factor loading without changing the observed distribution. Economic theory
must therefore decide whether the common factor moves $Y_2$ and $Y_1$ in the
same direction or in opposite directions.

\subsection{Cumulant notation}

Let $\Phi_{W_2,W_1}$ be the log characteristic function of
$(W_{2,t+1},W_{1,t+1})$, and let $\Phi_{W_2}$ be the log characteristic function of
$W_{2,t+1}$. Define joint cumulants by
\begin{equation}
\Phi_{W_2,W_1}^{k,l}:=
\left.
\frac{\partial^{k+l}\Phi_{W_2,W_1}(\zeta,\xi)}
     {i^{k+l}\partial\zeta^k\partial\xi^l}
\right|_{\zeta=0,\xi=0},
\qquad
\Phi_{W_2}^{k}:=
\left.
\frac{\partial^{k}\Phi_{W_2}(\zeta)}
     {i^{k}\partial\zeta^k}
\right|_{\zeta=0}.
\label{eq:cumulants}
\end{equation}
For ordinary intuition, $\Phi_{W_2}^2$ is the variance of $W_2$, $\Phi_{W_2}^3$
is its third cumulant, and higher orders measure higher-order departures from
Gaussianity. Gaussian variables have all cumulants above order two equal to zero.

Because cumulants of independent sums add, the latent model implies, for any
nonnegative integer $p$,
\begin{align}
\Phi_{W_2}^{3+p} &= \kappa_u^{3+p}+\kappa_v^{3+p},\label{eq:kappa1}\\
\Phi_{W_2,W_1}^{2+p,1} &= \alpha\kappa_u^{3+p}+\theta\kappa_v^{3+p},\label{eq:kappa2}\\
\Phi_{W_2,W_1}^{1+p,2} &= \alpha^2\kappa_u^{3+p}+\theta^2\kappa_v^{3+p},
\label{eq:kappa3}
\end{align}
where $\kappa_u^{m}$ and $\kappa_v^{m}$ are the order-$m$ cumulants of $u$ and
$v$. Eliminating these nuisance cumulants gives the observable moment equation
\begin{equation}
M_p(\alpha,\theta):=
\Phi_{W_2,W_1}^{1+p,2}-\alpha^2\Phi_{W_2}^{3+p}
-(\theta+\alpha)\left(\Phi_{W_2,W_1}^{2+p,1}-\alpha\Phi_{W_2}^{3+p}\right)=0.
\label{eq:Mp}
\end{equation}
This is the scalar version of the moment condition implemented by \texttt{trigmm}.

\subsection{Identification from two cumulant equations}

Two different values of $p$ give two equations in the two unknowns $(\alpha,\theta)$.
For example,
\begin{equation}
M_0(\alpha,\theta)=0,
\qquad
M_1(\alpha,\theta)=0
\label{eq:M01}
\end{equation}
can point identify $(\alpha,\theta)$. The default \texttt{trigmm} specification
uses $p=0,1$. The command can also use $p=0,1,2$, which gives three moment
equations for two parameters and therefore overidentifying restrictions.

For a pair $q<\widetilde q$, the essential nondegeneracy condition is that the
higher cumulants are not lined up in a way that makes the two equations carry the
same information. In the notation here, a convenient statement is
\begin{equation}
\Phi_{W_2}^{3+\widetilde q}\,\Phi_{W_2,W_1}^{2+q,1}
\ne
\Phi_{W_2}^{3+q}\,\Phi_{W_2,W_1}^{2+\widetilde q,1}.
\label{eq:nondegenerate}
\end{equation}
For the default pair $q=0$, $\widetilde q=1$, this condition can fail if either
$u$ or $v$ is Gaussian, if both are symmetric, or if they have exactly the same
distribution. This is why the method is described as identification by
non-Gaussianity: information must be present beyond means and variances.

Once $(\alpha,\theta)$ is identified, the loading of the common factor is
\begin{equation}
\beta_u=\alpha-\theta.
\label{eq:beta-u}
\end{equation}
The latent variances can be recovered from second moments. Let
$\sigma_u^2:=\Var(u)$, $\sigma_v^2:=\Var(v)$, and $\sigma_r^2:=\Var(r)$. Then
\begin{align}
\sigma_u^2 &=
\frac{\Cov(W_1,W_2)-\theta\Var(W_2)}{\alpha-\theta},
\label{eq:varu}\\
\sigma_v^2 &= \Var(W_2)-\sigma_u^2,
\label{eq:varv}\\
\sigma_r^2 &= \Var(W_1)-\alpha^2\sigma_u^2-\theta^2\sigma_v^2.
\label{eq:varr}
\end{align}
The nonnegativity and plausible sizes of these implied variances are useful
model diagnostics.

\subsection{How one scalar external instrument can enter}

The internal \texttt{trigmm} moments do not require an external instrument. If a
scalar external instrument $Z_{t+1}$ is available and satisfies the conventional
exclusion condition
\begin{equation}
\E\left[Z_{t+1}\{W_{1,t+1}-\theta W_{2,t+1}\}\right]=0,
\label{eq:external-iv}
\end{equation}
then it can be added as an additional GMM moment. With one external instrument
and the internal moment pair $M_0=M_1=0$, the model is overidentified. This is a
particularly useful implementation because it lets the researcher compare three
sources of information:
\begin{enumerate}
\item the conventional IV estimate from \eqref{eq:external-iv};
\item the internal higher-cumulant estimate from \eqref{eq:Mp};
\item the combined GMM estimate using both.
\end{enumerate}
Agreement across these estimates is not proof, but it is a strong robustness
check because the identifying assumptions are very different.

\subsection{Inference and computation}

The point-identified \texttt{trigmm} estimator is a nonlinear GMM estimator. Under
strong identification and valid moments, standard GMM covariance formulas give
standard errors for $\theta$, $\alpha$, $\beta_u$, the variance components, and the
recovered $X_t$ coefficients. The paper emphasizes three practical points.
\begin{itemize}
\item The objective function can have multiple local minima. Starting values
matter, especially with the overidentified $p=0,1,2$ specification.
\item The overidentified specification allows a Hansen $J$ test. When the extra
moments come from higher cumulants, a rejection indicates a failure of the
latent distributional assumptions or the linear model, not merely a bad external
instrument.
\item Weak or nearly degenerate higher-moment identification shows up as large
standard errors and poorly conditioned moment derivatives. The paper suggests
checking singular values of the GMM derivative matrix as a diagnostic.
\end{itemize}

\section{Direct option 2 from Lee--Lewbel--Schennach--Zhang: second-moment set identification}
\label{sec:lsz-set}

The same common-factor model also gives a set-identified result that uses only
second moments. This option is useful when the non-Gaussianity needed for
Section~\ref{sec:lsz-point} is not credible or is too weak to estimate precisely.

The underlying Lee--Lewbel--Schennach--Zhang theory contains a sharp identified
set based on decomposing the observed distribution into independent Gaussian and
non-Gaussian components. That sharp set is conceptually the most complete
second-moment object, but it is difficult to estimate and explain in applied
work. The Stata paper therefore focuses on a coarser but transparent outer bound
based only on the covariance matrix of $(W_2,W_1)$. The formulas below describe
that implemented bound.

Assume first that $\beta_u=\alpha-\theta>0$. Define
\begin{equation}
B_0:=\frac{\E[W_{1,t+1}W_{2,t+1}]}{\E[W_{2,t+1}^2]},
\qquad
D_0:=
\frac{\E[W_{1,t+1}^2]\E[W_{2,t+1}^2]-\{\E[W_{1,t+1}W_{2,t+1}]\}^2}
     {\{\E[W_{2,t+1}^2]\}^2}.
\label{eq:B0D0}
\end{equation}
If the residuals are mean zero, $B_0$ is simply the population OLS slope from
regressing $W_1$ on $W_2$, and $D_0$ is the residual variance from that regression
scaled by $\Var(W_2)$.

Because
\begin{equation}
B_0=
\frac{\alpha\E[u^2]+\theta\E[v^2]}{\E[u^2]+\E[v^2]}
=\alpha\lambda+\theta(1-\lambda),
\qquad
\lambda:=\frac{\E[u^2]}{\E[u^2]+\E[v^2]}\in[0,1],
\label{eq:B0-convex}
\end{equation}
we obtain the first bound,
\begin{equation}
\theta\le B_0\le\alpha.
\label{eq:lsz-order}
\end{equation}
The sharper second-moment bound reported in the paper is
\begin{equation}
(B_0-\theta)(\alpha-B_0)\le D_0.
\label{eq:lsz-hyperbola}
\end{equation}
Thus the population identified set is
\begin{equation}
\Theta_{\text{LSZ,set}}^+ :=
\{(\alpha,\theta):\theta\le B_0\le\alpha,
\ (B_0-\theta)(\alpha-B_0)\le D_0\}.
\label{eq:lsz-set-positive}
\end{equation}
The plus sign indicates $\beta_u>0$.

If instead $\beta_u<0$, apply the positive-loading result to $-W_2$. The order
reverses and the set becomes
\begin{equation}
\Theta_{\text{LSZ,set}}^- :=
\{(\alpha,\theta):\alpha\le B_0\le\theta,
\ (\theta-B_0)(B_0-\alpha)\le D_0\}.
\label{eq:lsz-set-negative}
\end{equation}

\subsection{Important consequence: the marginal set for $\theta$ is usually not a finite interval}

The joint set \eqref{eq:lsz-set-positive} is informative, but its projection onto
$\theta$ alone is
\begin{equation}
\{\theta:\exists\alpha\text{ such that }(\alpha,\theta)\in
\Theta_{\text{LSZ,set}}^+\}=(-\infty,B_0]
\label{eq:theta-halfline-plus}
\end{equation}
in population. For $\beta_u<0$, the projection is $[B_0,\infty)$. Therefore the
second-moment set is a joint bound on the causal effect $\theta$ and the total
common-factor effect $\alpha$, not a finite scalar interval for $\theta$ by itself.
A finite interval for $\theta$ requires one more substantive restriction, such as a
known range for $\beta_u=\alpha-\theta$ or a known range for $\alpha$.

For example, if $\beta_u>0$ and $0\le \beta_u\le\bar\beta$, then
$\alpha=\theta+\beta_u$ adds
\begin{equation}
0\le \alpha-\theta\le\bar\beta.
\label{eq:beta-bound}
\end{equation}
Projecting the intersection of \eqref{eq:lsz-set-positive} and
\eqref{eq:beta-bound} gives a finite interval whenever $\bar\beta$ is finite.
This extra magnitude restriction is not part of the base \texttt{trigmmset}
second-moment result; it must come from outside the paper.

\subsection{Confidence region used by \texttt{trigmmset}}

In sample, estimate $B_0$ and $D_0$ by plug-in estimators $(\widehat B,\widehat D)$
and estimate their covariance matrix $\widehat V$ by the delta method. For a
significance level $\delta$, construct the ellipse
\begin{equation}
\cE_\delta:=
\left\{(B,D):
\begin{bmatrix}B-\widehat B & D-\widehat D\end{bmatrix}
\widehat V^{-1}
\begin{bmatrix}B-\widehat B\\ D-\widehat D\end{bmatrix}
\le \chi^2_{2,1-\delta}
\right\}.
\label{eq:ellipse}
\end{equation}
Then $\cE_\delta$ contains the true $(B_0,D_0)$ with approximate probability
$1-\delta$. The command reports a worst-case region that contains the union of
all identified sets generated by all $(B,D)\in\cE_\delta$. For $\beta_u>0$, define
\begin{align}
B^-&:=\min_{(B,D)\in\cE_\delta}B,
&
B^+&:=\max_{(B,D)\in\cE_\delta}B,
\label{eq:Bminusplus}\\
D^*&:=\max_{(B,D)\in\cE_\delta}
\{D+(B-B^+)(B^- -B)\}.
\label{eq:Dstar}
\end{align}
The reported confidence set is the union of four simple regions:
\begin{align}
&\alpha\ge\theta,\quad B^-\le\alpha\le B^+,
\quad B^-\le\theta\le B^+,\label{eq:region1}\\
&B^-\le\alpha\le B^+,
\quad \theta\le B^-,\label{eq:region2}\\
&B^-\le\theta\le B^+,
\quad B^+\le\alpha,\label{eq:region3}\\
&D^*\ge(\alpha-B^+)(B^- -\theta),
\quad B^+\le\alpha,
\quad \theta\le B^-.
\label{eq:region4}
\end{align}
This construction is conservative but easy to draw and easy to report. It gives a
confidence region for the joint object $(\alpha,\theta)$. The marginal projection
onto $\theta$ is again generally one-sided: approximately $(-\infty,B^+]$ for
$\beta_u>0$ and $[B^-,\infty)$ for $\beta_u<0$.

\subsection{Coefficients on controls}

The endogenous-equation slope $\beta_2^R$ is point identified by the projection of
$Y_2$ on $X_t$. The coefficient vector on $X_t$ in the outcome equation is
\begin{equation}
\beta_1(\theta)=\beta_1^R-\beta_2^R\theta.
\label{eq:lsz-betarecovery}
\end{equation}
Because $\theta$ is set identified, $\beta_1$ is also set identified. The
\texttt{trigmmset} paper describes a conservative element-by-element way to bound
these coefficients using confidence intervals for the components of $\beta_2^R$,
$\beta_1^R$, and $B_0$. In the scalar notation here, if the sign of a particular
component of $\beta_2^R$ is known, the one-sided bound on $\theta$ translates into a
one-sided bound on that element of $\beta_1$. If the sign of that component is not
statistically determined, the corresponding coefficient bound may be
uninformative.

\section{Options from Lewis's review, translated to the scalar triangular setting}
\label{sec:lewis-options}

Lewis's Annual Review paper is written for structural vector autoregressions and
simultaneous systems with a square impact matrix. It is therefore not identical
to the Lewbel triangular model in \eqref{eq:structural}--\eqref{eq:firststage},
and it is also not identical to the Lee--Lewbel--Schennach--Zhang three-shock
common-factor model. To use those options in the present scalar problem, impose
an alternative square decomposition of the two residuals:
\begin{equation}
U_{t+1}:=\begin{bmatrix}W_{2,t+1}\\ W_{1,t+1}\end{bmatrix}
=B\eta_{t+1},
\qquad
B=\begin{bmatrix}b_{11}&b_{12}\\ b_{21}&b_{22}\end{bmatrix}.
\label{eq:square-B}
\end{equation}
The shocks $\eta_{t+1}=(\eta_{1,t+1},\eta_{2,t+1})^\T$ are normalized in some
way, such as unit variances or a unit diagonal. After estimating $B$, the scalar
coefficient analogous to $\theta$ is an impact ratio for a labeled shock. If
column $j$ is interpreted as the shock that moves $W_2$, then
\begin{equation}
\theta_j:=\frac{b_{2j}}{b_{1j}},
\qquad b_{1j}\ne0.
\label{eq:impact-ratio}
\end{equation}
This ratio is not automatically the structural $\theta$ in
\eqref{eq:resid-system}; it becomes the target only if the researcher is willing
to interpret the selected shock as the exogenous variation in $W_2$ whose causal
impact on $W_1$ is sought. Shock labeling is therefore central.

\subsection{Heteroskedasticity option A: two variance regimes}
\label{subsec:regimes}

Let the scalar variable $Z_{t+1}$ partition the data into two regimes, say
$Z_{t+1}=0$ and $Z_{t+1}=1$. Suppose
\begin{equation}
\Sigma_r:=\Var(U_{t+1}\mid Z_{t+1}=r)=B\Lambda_rB^\T,
\qquad
\Lambda_r=\diag(\lambda_{r1},\lambda_{r2}),
\qquad r=0,1.
\label{eq:regime-covs}
\end{equation}
The impact matrix $B$ is constant across regimes, but the shock variances change.
The two-regime heteroskedasticity condition is
\begin{equation}
\frac{\lambda_{11}}{\lambda_{01}}
\ne
\frac{\lambda_{12}}{\lambda_{02}}.
\label{eq:nonprop}
\end{equation}
This says that the two shock variances do not change proportionally across the
two regimes. Then
\begin{equation}
\Sigma_1\Sigma_0^{-1}=B\Lambda_1\Lambda_0^{-1}B^{-1},
\label{eq:eigen}
\end{equation}
so the columns of $B$ are eigenvectors of an observable matrix, up to scale and
order. Once a column is labeled as the $W_2$ shock, $\theta_j=b_{2j}/b_{1j}$ is
identified.

A particularly transparent special case is that only the variance of the labeled
shock $j$ changes. Then
\begin{equation}
\Sigma_1-\Sigma_0=(\lambda_{1j}-\lambda_{0j})b_jb_j^\T,
\label{eq:onediff}
\end{equation}
where $b_j$ is the $j$-th column of $B$. With
$U=(W_2,W_1)^\T$, the impact ratio is
\begin{equation}
\theta_j=
\frac{\Cov(W_2,W_1\mid Z=1)-\Cov(W_2,W_1\mid Z=0)}
     {\Var(W_2\mid Z=1)-\Var(W_2\mid Z=0)}.
\label{eq:rigobon-ratio}
\end{equation}
This is the closest variance-regime analogue of a single-instrument estimator.
The scalar $Z$ is not an excluded IV; it supplies a high-variance and a low-variance
state.

\subsection{Heteroskedasticity option B: more regimes, Markov switching, smooth transition, and GARCH}

The same covariance-equation logic extends in several directions.
\begin{itemize}
\item \textbf{More than two regimes.} Additional regimes add more covariance
matrices. If $B$ is constant and the diagonal variance changes are rich enough,
$B$ is overidentified. Overidentifying restrictions can test whether the fixed
impact matrix and diagonal shock-variance assumptions are plausible.
\item \textbf{Markov-switching variances.} The regimes are latent rather than
observed. Identification still comes from distinct variance matrices satisfying
a nonproportionality condition, but estimation must also infer the regimes.
\item \textbf{Smooth transition variances.} The covariance matrix evolves as
\begin{equation}
\Omega_t=(1-F(s_t))\Sigma_0+F(s_t)\Sigma_1,
\label{eq:smooth}
\end{equation}
where $s_t$ is a transition variable. Conditional on the transition parameters,
the endpoints $\Sigma_0$ and $\Sigma_1$ identify $B$ by the same logic as
Section~\ref{subsec:regimes}. Identification is weak if the transition is nearly
flat or if $s_t$ does not track actual volatility changes.
\item \textbf{GARCH or stochastic volatility.} Parametric volatility models use
the full time path of conditional variances. The maintained idea remains fixed
$B$ plus diagonal, time-varying structural variances. Parametric assumptions can
improve efficiency but introduce functional-form risk.
\end{itemize}
In all of these variants, one can recover an impact ratio such as
\eqref{eq:impact-ratio} only after normalizing and labeling the relevant column
of $B$.

\subsection{Heteroskedasticity option C: unconditional moments of squared residuals}

Lewis also reviews distribution-free identification based on persistence in
volatility. In the present two-variable notation, the useful objects are
autocovariances of squared and cross-product residuals, such as
\begin{equation}
\Cov\left(\operatorname{vec}(U_{t+1}U_{t+1}^\T),
\operatorname{vec}(U_{t+1-s}U_{t+1-s}^\T)^\T\right),
\qquad s>0.
\label{eq:squared-autocovs}
\end{equation}
If structural shocks themselves are serially uncorrelated but their variances are
persistent, these autocovariances reveal how the structural variances move over
time. Identification requires the volatility processes of the shocks not to have
proportional persistence patterns. This option is attractive because it does not
require choosing exact regimes, but it is demanding in finite samples because it
uses fourth-order quantities and their autocovariances.

\subsection{Non-Gaussianity option A: independent component analysis}

Instead of changing variances over time, assume that the components of
$\eta_{t+1}$ in \eqref{eq:square-B} are mutually independent and that at most one
of them is Gaussian. Then, by the Darmois--Skitovich/Comon logic reviewed by
Lewis, the decomposition $U=B\eta$ is unique up to scale, sign, and column order.
In the two-shock case this means that at least one shock must be non-Gaussian.

Independent component analysis (ICA) estimates the rotation of the whitened
residual vector that makes the candidate shocks as independent or as
non-Gaussian as possible. After estimating $B$, the scalar causal parameter is
again an impact ratio $b_{2j}/b_{1j}$ for the economically labeled column. The
single scalar $Z$ can help label the column, for example by choosing the shock
most correlated with $Z$ or the shock whose response signs match prior theory,
but the statistical identification itself comes from independence and
non-Gaussianity, not from the exclusion restriction $\E[Z\varepsilon_1]=0$.

\subsection{Non-Gaussianity option B: likelihood and pseudo-likelihood}

A parametric likelihood version specifies densities for the independent shocks,
for example Student-$t$ densities. If the density is correctly specified,
maximum likelihood can be efficient. Pseudo-maximum likelihood can remain
consistent under certain misspecifications, provided the pseudo-model is itself
identified. The cost is that there are many possible ways to be non-Gaussian, so
a particular density choice is an additional modeling decision.

In the scalar triangular report, this option again delivers $\theta$ only through
an estimated square mixing matrix and a labeled impact ratio. It is not the same
as the Lee--Lewbel--Schennach--Zhang cumulant equations, which are tailored to
the three-component common-factor model.

\subsection{Non-Gaussianity option C: moment-based GMM with coskewness or cokurtosis}

Moment-based non-Gaussian identification selects particular independence-implied
moments. Let $\eta=C U$ denote candidate shocks after applying an unmixing matrix
$C=B^{-1}$ and normalizing variances to one. Typical restrictions include
\begin{align}
\E[\eta_1\eta_2]&=0,
\label{eq:orthog}\\
\E[\eta_1^2\eta_2]&=0,
\qquad
\E[\eta_1\eta_2^2]=0,
\label{eq:coskew}\\
\E[\eta_1^2\eta_2^2]-1&=0,
\qquad
\E[\eta_1^3\eta_2]=0,
\qquad
\E[\eta_1\eta_2^3]=0.
\label{eq:cokurt}
\end{align}
The restrictions in \eqref{eq:coskew} use skewness-type information; those in
\eqref{eq:cokurt} use kurtosis-type information. Identification requires at
least one shock to have relevant nonzero skewness or nonzero excess kurtosis. A
GMM estimator chooses $C$ so these cross-moments are close to zero, then obtains
$B=C^{-1}$ and the impact ratio \eqref{eq:impact-ratio}.

This option is transparent because the identifying moments are explicit, but it
can be noisy: estimating fourth moments and the covariance matrix of moment
conditions may require moments up to order eight.

\subsection{Non-independent components and coheteroskedasticity}

The Annual Review also discusses newer work relaxing full independence. The core
idea is to replace full independence by restrictions on higher-order cumulant
tensors, such as diagonal cumulant tensors or reflectional-invariance patterns.
These methods are promising because they can allow some forms of dependence,
including coheteroskedasticity. In the scalar triangular problem they are best
viewed as advanced alternatives to the simple square ICA model, not as direct
Lewbel or \texttt{trigmm} estimators.

\section{Using one scalar instrument with each option}
\label{sec:singleZ}

The same observed scalar $Z$ can play different roles. Confusing these roles is
the easiest way to misuse these methods.

\begin{enumerate}[label=\textbf{Role \arabic*.}]
\item \textbf{External IV.} If $\E[Z\varepsilon_1]=0$ and $\E[ZW_2]\ne0$, then
$Z$ directly identifies $\theta$ by conventional IV:
\begin{equation}
\theta_{\text{IV}}=\frac{\E[ZW_1]}{\E[ZW_2]}.
\label{eq:directiv}
\end{equation}
This moment can be combined with \texttt{trigmm}'s internal moments.

\item \textbf{Lewbel heteroskedasticity shifter.} If $Z$ satisfies
\eqref{eq:lewbel-exact-assump}, then $(Z-\E Z)W_2$ is the constructed
instrument. Here the relevant covariance is $\Cov(Z,W_2^2)$, not $\E[ZW_2]$.

\item \textbf{Variance-regime indicator or volatility state.} If $Z$ defines
high- and low-volatility regimes, then it supplies two covariance matrices. The
assumption is not exclusion but fixed $B$ and nonproportional changes in shock
variances.

\item \textbf{Shock-labeling variable.} In ICA or non-Gaussian square models, $Z$
may help decide which column of $B$ is economically meaningful. It does not by
itself identify the rotation unless an additional IV or moment condition is
imposed.

\item \textbf{No role in internal LSZ identification.} The Lee--Lewbel--Schennach--Zhang
point and set identification results do not require $Z$. They can be used when no
external instrument exists. A valid $Z$ can be added, but it is optional.
\end{enumerate}

\section{Inference and weak-identification issues}
\label{sec:inference}

All the methods in this report use information beyond the first two unconditional
moments of a homoskedastic linear regression. That information can be weak.

\heading{Lewbel exact and bounded-invalidity inference.}
If $\Cov(Z,W_2^2)$ is near zero, the constructed instrument is weak. Standard IV
weak-instrument diagnostics are informative because the estimator is implemented
as two-stage least squares with the constructed instrument. In the bounded
invalidity interval, the same weak relevance makes the interval wide and can make
sample roots unstable.

\heading{Lee--Lewbel--Schennach--Zhang point inference.}
If the relevant higher cumulants of $u$ and $v$ are small or nearly degenerate,
the GMM objective is flat in some direction. Then standard errors can be large,
local minima become a practical problem, and Wald intervals can be misleading.
The overidentified $p=0,1,2$ specification, the Hansen $J$ test, and singular
values of the moment derivative matrix are useful diagnostics, but none proves
identification strength in finite samples.

\heading{Lee--Lewbel--Schennach--Zhang set inference.}
The \texttt{trigmmset} confidence region is explicitly set-valued. It first builds
an ellipse for $(B_0,D_0)$ and then reports a conservative worst-case joint
region for $(\alpha,\theta)$. This is more honest than reporting a spurious
finite scalar confidence interval for $\theta$ when the population identified set
itself projects to a half-line.

\heading{Variance-regime inference.}
If variances change little, or if both shock variances change almost
proportionally, the eigenvectors in \eqref{eq:eigen} are weakly identified. Robust
GMM confidence sets or weak-identification-robust procedures are preferable to
simple Wald intervals.

\heading{Non-Gaussian inference.}
If the shocks are close to Gaussian, independence-based rotations are weakly
identified. Testing for non-Gaussianity and independence is useful but difficult
because the shocks are estimated objects. Robust score/projection methods have
been developed in the literature, but they are more complex than ordinary
standard errors.

\section{Software and implementation notes}
\label{sec:software}

\begin{itemize}
\item \textbf{Stata.} Lee et al. provide \texttt{trigmm} for the point-identified
non-Gaussian common-factor method and \texttt{trigmmset} for the second-moment
set. The commands are wrappers around Stata's \texttt{gmm} infrastructure and use
Stata post-estimation tools. Lewbel's heteroskedasticity-generated instruments
are implemented in \texttt{ivreg2h}.

\item \textbf{R or Python.} The scalar formulas in Sections~\ref{sec:lewbelbaseline}
through \ref{sec:lsz-set} are simple enough to implement directly. For the LSZ
point estimator, one needs routines for sample cumulants and nonlinear GMM. For
set estimation, one needs only plug-in estimates of $B_0$ and $D_0$, their
delta-method covariance matrix, and a routine to draw or optimize over the
reported joint region. Generic GMM, nonlinear optimization, and bootstrap tools
are sufficient, but a direct off-the-shelf implementation of \texttt{trigmm} for R
is not part of the attached papers.

\item \textbf{ICA and SVAR methods.} Generic ICA, VAR, GARCH, and structural-VAR
software can estimate some of the broader Lewis-type models. The researcher must
still map the estimated impact matrix into the scalar ratio \eqref{eq:impact-ratio}
and justify the shock label. These packages do not automatically estimate the
Lewbel triangular coefficient $\theta$.
\end{itemize}

\section{A practical decision tree}
\label{sec:decision}

\begin{enumerate}
\item \textbf{Start with the maintained triangular model.} Residualize both
$Y_{1,t+1}$ and $Y_{2,t+1}$ on the same $X_t$ and work with
$W_{1,t+1},W_{2,t+1}$.

\item \textbf{If a conventional scalar IV is credible, use it and report it.}
Estimate \eqref{eq:directiv} or its controlled version. Then consider adding
Lewbel or \texttt{trigmm} moments as robustness checks or efficiency-improving
overidentifying restrictions.

\item \textbf{If no conventional IV is credible but a heteroskedasticity shifter is
credible, use Lewbel.} Check whether $\Cov(Z,W_2^2)$ is meaningfully nonzero. If
exact validity is doubtful, report the bounded-invalidity interval over a grid of
$\tau$ values.

\item \textbf{If the common-factor story is credible and the errors are plausibly
non-Gaussian, use \texttt{trigmm}.} Specify the sign of $\beta_u$, estimate with
$p=0,1$ and with $p=0,1,2$, vary starting values, and report diagnostics.

\item \textbf{If the common-factor story is credible but non-Gaussianity is not,
use \texttt{trigmmset}.} Report the joint region for $(\alpha,\theta)$ and state
clearly whether the implied marginal set for $\theta$ is one-sided. Do not report
a finite interval for $\theta$ unless an additional magnitude restriction is
imposed.

\item \textbf{If the data are time series with credible volatility regimes or
persistent volatility, consider the Lewis heteroskedasticity options.} Make clear
that this replaces the common-factor model by a square impact-matrix model and
that $\theta$ is an impact ratio for a labeled shock.

\item \textbf{If non-Gaussian independent shocks are credible, consider ICA,
likelihood, or higher-moment GMM.} Again, report how the shock was labeled and
how the impact ratio was normalized.
\end{enumerate}

\section{Conclusion}
\label{sec:conclusion}

In the scalar triangular setting, the methods differ mainly in what replaces the
missing excluded instrument. Lewbel uses a scalar heteroskedasticity shifter to
construct an instrument. Lee--Lewbel--Schennach--Zhang use a common-factor error
structure: higher cumulants yield point identification, while second moments
yield a conservative joint set. Lewis's review offers a menu of broader
higher-moment strategies--variance regimes, volatility processes, ICA,
likelihood, and moment-based non-Gaussian GMM--but those strategies require a
square shock decomposition and a shock-labeling step before they deliver a
scalar effect comparable to $\theta$.

The most transparent reporting strategy is to state, for each estimate or set,
which object identifies $\theta$: an exclusion moment, a heteroskedasticity
covariance, a non-Gaussian cumulant system, a second-moment common-factor
inequality, a variance-regime covariance change, or an independent-component
rotation. Once that distinction is clear, all other coefficients follow from the
simple recovery formula $\beta_1(w)=\beta_1^R-\beta_2^Rw$.



\clearpage
\section{Inference and R bootstrap implementation guide}
\label{sec:r-bootstrap-guide}

The preceding sections describe what identifies the coefficient of the scalar
endogenous regressor.  This section describes how to attach standard errors,
confidence intervals, or confidence sets to the resulting estimate or identified
set.  The emphasis is deliberately practical: what can be done directly in R,
what requires a short custom wrapper, and what is currently easier to do in Stata
or by translating published Stata/MATLAB code.

There are three distinct inferential targets.  They should not be mixed.
\begin{enumerate}
    \item \emph{Point-identified inference}: report a standard error or
    confidence interval for a singleton estimate $\hat\theta$.
    \item \emph{Endpoint inference}: report uncertainty about the lower and
    upper endpoints $(L_0,U_0)$ of a scalar identified interval
    $\Theta_0=[L_0,U_0]$.
    \item \emph{Partial-identification confidence sets}: report a set that
    covers either the entire identified set $\Theta_0$ or the true unknown
    scalar value $\theta_0\in\Theta_0$.  These are different coverage goals.
\end{enumerate}

In the scalar triangular setting, the full model often reduces to a one-dimensional
coefficient set for $\theta$.  Once a candidate $w$ is retained, the other
coefficients are recovered from
\begin{equation}
\beta_1(w)=\beta_1^R-\beta_2^R w.
\end{equation}
Thus the most transparent workflow is: first form an estimate or confidence set
for $\theta$, and then map each endpoint or retained grid value through the
coefficient recovery formula.

\subsection{What was found in R}
\label{subsec:r-packages-found}

Table \ref{tab:r-packages} summarizes the R packages and code bases that are
most relevant to the current problem.  The important conclusion is that R has
usable point-estimation and bootstrap building blocks, but I did not find a
standard R package that directly implements the Lewbel (2012) Theorem 3
set-identified interval or the Lee--Lewbel--Schennach--Zhang \texttt{trigmmset}
confidence region.

\begin{longtable}{>{\raggedright\arraybackslash}p{0.20\textwidth}>
                  {\raggedright\arraybackslash}p{0.23\textwidth}>
                  {\raggedright\arraybackslash}p{0.29\textwidth}>
                  {\raggedright\arraybackslash}p{0.18\textwidth}}
\caption{R packages and code bases relevant to bootstrap inference in the scalar triangular model.}\label{tab:r-packages}\\
\toprule
Package or code base & Role in this problem & What it provides & Limitation \\
\midrule
\endfirsthead
\toprule
Package or code base & Role in this problem & What it provides & Limitation \\
\midrule
\endhead
\texttt{REndo::}\newline\texttt{hetErrorsIV()} & Closest direct R implementation of the Lewbel heteroskedastic-error triangular estimator. & Estimates linear models with endogenous regressors using heteroskedastic covariance restrictions; reports robust and cluster-adjusted standard errors. & Point-identified estimator only; no direct bootstrap argument for the set-identified bounds. \\
\addlinespace
\texttt{ivlewbel::}\newline\texttt{lewbel()} & Older direct Lewbel-style R implementation. & Implements heteroskedasticity-based GMM/IV estimation for triangular models. & Older and less maintained; point-estimation focus. \\
\addlinespace
\texttt{ivreg} & General IV/2SLS infrastructure. & After manually constructing Lewbel instruments, runs 2SLS and exposes an object compatible with many covariance and bootstrap tools. & Does not construct Lewbel instruments or set bounds automatically. \\
\addlinespace
\texttt{gmm} & User-supplied moment conditions. & Natural framework for hand-coded nonlinear \texttt{trigmm}-style moments, Lewbel moments, and custom GMM covariance estimation. & Bootstrap and set-inversion must be wrapped by the user. \\
\addlinespace
\texttt{boot} & General full-procedure bootstrap. & Resample rows, clusters, or blocks; return any statistic, including a point estimate, a lower endpoint, an upper endpoint, or a whole vector of moments. & Validity depends on the statistic; ordinary $n$-out-of-$n$ bootstrap can fail near nonsmooth set boundaries. \\
\addlinespace
\texttt{sandwich::}\newline\texttt{vcovBS()} & Generic bootstrap covariance for model objects. & Case or cluster bootstrap covariance matrices for supported objects through \texttt{update()} and \texttt{coef()}. & Less convenient for custom endpoint sets that are not model objects. \\
\addlinespace
\texttt{fwild}\newline\texttt{clusterboot} & Wild cluster bootstrap for linear and IV models. & Supports \texttt{ivreg} objects, so it can be used after manual construction of Lewbel instruments. & Point-identified IV inference; not a partial-identification set package. \\
\addlinespace
\texttt{clusterSEs::}\newline\texttt{cluster.wild}\newline\texttt{.ivreg()} & Alternative wild cluster bootstrap for IV. & Wild-cluster-bootstrap $p$-values and intervals for \texttt{ivreg} models. & Point-identified IV inference only. \\
\addlinespace
\texttt{ivmodel} & Weak-IV robust inference with one endogenous regressor. & Anderson--Rubin, conditional likelihood ratio, and related intervals for linear IV. & Designed for conventional linear IV, not for \texttt{trigmmset} or Lewbel Theorem 3 set endpoints. \\
\addlinespace
\texttt{moonboot} & $m$-out-of-$n$ bootstrap building block. & Useful when ordinary bootstrap is suspect for nonsmooth or irregular statistics. & General method; no triangular-system estimator. \\
\addlinespace
\texttt{ivmte::}\newline\texttt{boundCI()} & Useful R example of bootstrap inference for partial-identification bounds. & Takes original bounds and bootstrapped bounds and constructs confidence intervals for bounds. & For marginal treatment effect bounds, not for the scalar triangular model; useful mainly as a template. \\
\bottomrule
\end{longtable}

The practical R stack is therefore:
\begin{equation*}
\texttt{REndo}\text{ or }\texttt{ivreg}\text{ for point estimation},\qquad
\texttt{boot}\text{ for full-procedure resampling},\qquad
\texttt{gmm}\text{ for custom moments}.
\end{equation*}
For clustered data, replace row resampling by cluster resampling or use a wild
cluster bootstrap after the generated instrument has been constructed.

\subsection{Point-identified Lewbel estimator: full-procedure bootstrap}
\label{subsec:point-lewbel-bootstrap}

For Lewbel's exact point-identified triangular estimator, the bootstrap must
repeat the entire generated-instrument construction.  It is not enough to hold
$\hat\varepsilon_2$ fixed across bootstrap draws.

For each bootstrap draw $b=1,\dots,B$:
\begin{enumerate}
    \item resample observations, clusters, or time blocks according to the
    sampling design;
    \item regress $Y_2$ on $X$ in the bootstrap sample and compute
    $\hat\varepsilon_{2,i}^{*(b)}$;
    \item construct $H_i^{*(b)}=(Z_i^{*(b)}-\bar Z^{*(b)})
    \hat\varepsilon_{2,i}^{*(b)}$;
    \item estimate the structural equation by IV using $X^{*(b)}$ and
    $H^{*(b)}$ as instruments;
    \item store $\hat\theta^{*(b)}$ and any recovered coefficient
    $\hat\beta_1^{*(b)}=\hat\beta_1^{R,*(b)}-
    \hat\beta_2^{R,*(b)}\hat\theta^{*(b)}$.
\end{enumerate}
A percentile interval is the interval between the empirical $\alpha/2$ and
$1-\alpha/2$ quantiles of the bootstrap draws.  A studentized interval stores
$(\hat\theta^{*(b)}-\hat\theta)/\widehat{\se}^{*(b)}$ and inverts the resulting
quantiles.  A basic interval reflects bootstrap endpoint deviations around the
original estimate.

A minimal R skeleton is:

\begin{verbatim}
# install.packages(c("ivreg", "boot"))
library(ivreg)
library(boot)

lewbel_point <- function(dat, y1, y2, X, Z) {
  f2 <- reformulate(X, response = y2)
  rf <- lm(f2, data = dat)
  dat$eps2_hat <- resid(rf)
  dat$H_lewbel <- (dat[[Z]] - mean(dat[[Z]], na.rm = TRUE)) * dat$eps2_hat

  rhs <- paste(c(X, y2), collapse = " + ")
  inst <- paste(c(X, "H_lewbel"), collapse = " + ")
  fiv <- as.formula(paste(y1, "~", rhs, "|", inst))
  fit <- ivreg(fiv, data = dat)
  unname(coef(fit)[y2])
}

boot_lewbel_point <- function(dat, y1, y2, X, Z, R = 1999, seed = 123) {
  set.seed(seed)
  stat <- function(d, idx) lewbel_point(d[idx, , drop = FALSE], y1, y2, X, Z)
  out <- boot(dat, statistic = stat, R = R)
  list(theta_hat = out$t0,
       percentile = boot.ci(out, type = "perc"),
       basic      = boot.ci(out, type = "basic"),
       boot_obj   = out)
}
\end{verbatim}

For clustered data, replace the observation index by a cluster index.  For a
wild cluster bootstrap, a convenient workflow is to construct the Lewbel
instrument in the original data, estimate the IV model with \texttt{ivreg}, and
then apply \texttt{fwildclusterboot::boottest()} to the coefficient of $Y_2$.
This is a point-identified IV bootstrap; it does not bootstrap the Lewbel
set-identified interval.

\subsection{Lewbel Theorem 3: bootstrapping a scalar set-identified interval}
\label{subsec:lewbel-theorem3-bootstrap}

Lewbel's Theorem 3 relaxes exact validity by allowing
$\Cov(Z,\varepsilon_1\varepsilon_2)$ to be nonzero but bounded.  In the scalar
case, suppose that for a user-chosen $\tau\in[0,1)$,
\begin{equation}
|\Corr(Z,\varepsilon_1\varepsilon_2)|
\le
\tau |\Corr(Z,
\varepsilon_2^2)|.
\label{eq:tau-assump-consolidated}
\end{equation}
Let $W_1$ and $W_2$ be the residuals from regressions of $Y_1$ and $Y_2$ on
$X$.  Define the sample analogues of
\begin{align*}
C_{12,Z}&=\Cov(W_1W_2,Z), & C_{22,Z}&=\Cov(W_2^2,Z),\\
V_{12}&=\Var(W_1W_2), & V_{22}&=\Var(W_2^2),\\
C_{12,22}&=\Cov(W_1W_2,W_2^2).
\end{align*}
The estimated endpoints of the Lewbel interval are the two real roots of
\begin{equation}
A\theta^2+B\theta+C=0,
\label{eq:lewbel-bound-quadratic-consolidated}
\end{equation}
where
\begin{align}
A &= 1-\tau^2,\label{eq:quadA}\\
B &= 2\left(\frac{C_{12,22}}{V_{22}}\tau^2
       -\frac{C_{12,Z}}{C_{22,Z}}\right),\label{eq:quadB}\\
C &= \left(\frac{C_{12,Z}}{C_{22,Z}}\right)^2
       -\frac{V_{12}}{V_{22}}\tau^2.\label{eq:quadC}
\end{align}
When $\tau=0$, the two roots collapse to the exact Lewbel point estimate
$C_{12,Z}/C_{22,Z}$.

The full-procedure bootstrap for this set is:
\begin{enumerate}
    \item compute the original interval $[\hat L,\hat U]$ from the original data;
    \item resample observations, clusters, or blocks;
    \item in each draw, recompute $W_1^*,W_2^*$, all covariance terms, and the
    two roots $[\hat L^*,\hat U^*]$;
    \item build endpoint intervals or a partial-identification confidence set
    from the empirical distribution of $(\hat L^*,\hat U^*)$.
\end{enumerate}

A minimal R implementation is:

\begin{verbatim}
pop_cov <- function(a, b) {
  mean((a - mean(a, na.rm = TRUE)) *
       (b - mean(b, na.rm = TRUE)), na.rm = TRUE)
}

rhs_formula <- function(y, X) {
  if (length(X) == 0) as.formula(paste(y, "~ 1")) else reformulate(X, response = y)
}

lewbel_set_bounds <- function(dat, y1, y2, X, Z, tau, eps = 1e-10) {
  stopifnot(tau >= 0, tau < 1)

  W1 <- resid(lm(rhs_formula(y1, X), data = dat))
  W2 <- resid(lm(rhs_formula(y2, X), data = dat))
  z  <- dat[[Z]]

  A12 <- W1 * W2
  A22 <- W2^2

  c12z  <- pop_cov(A12, z)
  c22z  <- pop_cov(A22, z)
  v12   <- pop_cov(A12, A12)
  v22   <- pop_cov(A22, A22)
  c1222 <- pop_cov(A12, A22)

  if (!is.finite(c22z) || abs(c22z) < eps ||
      !is.finite(v22) || v22 < eps) {
    return(c(L = NA_real_, U = NA_real_))
  }

  A <- 1 - tau^2
  B <- 2 * ((c1222 / v22) * tau^2 - (c12z / c22z))
  C <- (c12z / c22z)^2 - (v12 / v22) * tau^2

  disc <- B^2 - 4 * A * C
  if (!is.finite(disc) || disc < 0) return(c(L = NA_real_, U = NA_real_))

  roots <- sort((-B + c(-1, 1) * sqrt(disc)) / (2 * A))
  c(L = roots[1], U = roots[2])
}

bootstrap_lewbel_set <- function(dat, y1, y2, X, Z, tau,
                                 R = 1999, alpha = 0.05,
                                 m = nrow(dat), replace = TRUE,
                                 seed = 123) {
  set.seed(seed)
  n <- nrow(dat)

  theta_hat <- lewbel_set_bounds(dat, y1, y2, X, Z, tau)

  reps <- replicate(R, {
    ii <- sample.int(n, size = m, replace = replace)
    lewbel_set_bounds(dat[ii, , drop = FALSE], y1, y2, X, Z, tau)
  })

  reps <- t(reps)
  reps <- reps[complete.cases(reps), , drop = FALSE]

  list(
    point_bounds = theta_hat,
    percentile_outer = c(
      L = unname(quantile(reps[, "L"], alpha / 2, na.rm = TRUE)),
      U = unname(quantile(reps[, "U"], 1 - alpha / 2, na.rm = TRUE))
    ),
    basic_outer = c(
      L = 2 * theta_hat["L"] -
          unname(quantile(reps[, "L"], 1 - alpha / 2, na.rm = TRUE)),
      U = 2 * theta_hat["U"] -
          unname(quantile(reps[, "U"], alpha / 2, na.rm = TRUE))
    ),
    boot_reps = reps
  )
}
\end{verbatim}

The interval named \texttt{percentile\_outer} is easy to explain: it expands the
estimated lower endpoint downward and the estimated upper endpoint upward using
the empirical endpoint distributions.  This is a confidence interval for the
estimated set's endpoints, not automatically an optimal Imbens--Manski interval
for the unknown scalar $\theta_0\in[L_0,U_0]$.  For a paper, report which
coverage target is being used.

\subsection{Imbens--Manski/Stoye endpoint inference}
\label{subsec:im-stoye}

When the scalar identified set is $\Theta_0=[L_0,U_0]$ and the endpoints are
estimated by $(\hat L,\hat U)$, the most common partial-identification interval
has the form
\begin{equation}
\left[
\hat L-c_n \widehat{\se}(\hat L),\quad
\hat U+c_n \widehat{\se}(\hat U)
\right].
\label{eq:im-form}
\end{equation}
For set coverage, a conservative choice is the Bonferroni value $c_n=z_{1-
\alpha/2}$, or the projection of a joint bootstrap confidence region for
$(L_0,U_0)$.  For coverage of the true scalar point $\theta_0\in[L_0,U_0]$, the
Imbens--Manski critical value is generally smaller and depends on the estimated
width of the identified set relative to endpoint uncertainty.  Stoye's refinement
is important near point identification because naive endpoint intervals can have
nonuniform coverage when $U_0-L_0$ is close to zero.

A practical algorithm using the bootstrap draws from Section
\ref{subsec:lewbel-theorem3-bootstrap} is:
\begin{enumerate}
    \item estimate $[\hat L,\hat U]$ and bootstrap draws
    $[\hat L^{*(b)},\hat U^{*(b)}]$;
    \item estimate the covariance matrix of $(\hat L,\hat U)$ from the bootstrap
    draws;
    \item for a conservative interval, project the bootstrap joint region or use
    the Bonferroni outer interval;
    \item for an Imbens--Manski/Stoye interval, solve the published critical
    value equation using the endpoint standard errors and the estimated width
    $\hat U-\hat L$.
\end{enumerate}

This method is most appropriate for Lewbel Theorem 3, because that theorem gives
a scalar interval directly.  It is less natural for \texttt{trigmmset} unless the
joint $(\alpha,\gamma)$ set has first been projected to a finite scalar interval
for $\gamma$.

\subsection{A \texttt{trigmmset}-style second-moment bootstrap in R}
\label{subsec:trigmmset-bootstrap}

Lee--Lewbel--Schennach--Zhang's \texttt{trigmmset} command uses two observable
second-moment objects.  In the no-covariate residualized notation of this
report,
\begin{equation}
B_0=\frac{\E(W_1W_2)}{\E(W_2^2)},\qquad
D_0=\frac{\E(W_1^2)\E(W_2^2)-\E(W_1W_2)^2}{\E(W_2^2)^2}.
\label{eq:BD-consolidated}
\end{equation}
For positive common-factor loading, the structural parameters satisfy
\begin{equation}
\gamma\le B_0\le \alpha,
\qquad
(B_0-\gamma)(\alpha-B_0)\le D_0.
\label{eq:trigmmset-region-consolidated}
\end{equation}
The Stata command estimates $(B_0,D_0)$, uses a delta-method covariance matrix,
forms an ellipse in $(B,D)$ space, and reports a conservative worst-case region
for $(\alpha,\gamma)$ over that ellipse.

In R, a simple bootstrap building block is to resample $(\hat B,
\hat D)$:

\begin{verbatim}
trigmmset_BD <- function(dat, w, y) {
  ww <- dat[[w]]
  yy <- dat[[y]]

  Ey2 <- mean(yy^2, na.rm = TRUE)
  Ewy <- mean(ww * yy, na.rm = TRUE)
  Ew2 <- mean(ww^2, na.rm = TRUE)

  B <- Ewy / Ey2
  D <- (Ew2 * Ey2 - Ewy^2) / Ey2^2

  c(B = B, D = D)
}

bootstrap_trigmmset_BD <- function(dat, w, y, R = 1999,
                                   alpha = 0.05, seed = 123) {
  set.seed(seed)
  n <- nrow(dat)

  bd_hat <- trigmmset_BD(dat, w, y)

  reps <- replicate(R, {
    ii <- sample.int(n, n, replace = TRUE)
    trigmmset_BD(dat[ii, , drop = FALSE], w, y)
  })

  reps <- t(reps)

  list(
    BD_hat = bd_hat,
    BD_percentile_CI = apply(
      reps, 2, quantile,
      probs = c(alpha / 2, 1 - alpha / 2),
      na.rm = TRUE
    ),
    BD_cov = cov(reps, use = "complete.obs"),
    BD_reps = reps
  )
}
\end{verbatim}

This code only bootstraps $(B_0,D_0)$.  To reproduce the logic of
\texttt{trigmmset}, one must turn the joint uncertainty in $(B,D)$ into a joint
confidence region and then map that region through the inequalities in
\eqref{eq:trigmmset-region-consolidated}.  A bootstrap analogue of the Stata
ellipse is:
\begin{enumerate}
    \item compute $\hat\mu=(\hat B,\hat D)^\T$ and bootstrap draws
    $\hat\mu^{*(b)}=(\hat B^{*(b)},\hat D^{*(b)})^\T$;
    \item compute a covariance matrix $\hat V_{BD}$ from the bootstrap draws;
    \item form the ellipsoid
    \begin{equation}
    \mathcal E_{\alpha}^{boot}=
    \left\{\mu:
    (\mu-\hat\mu)^\T\hat V_{BD}^{-1}(\mu-\hat\mu)
    \le c_{1-\alpha}^{boot}
    \right\},
    \end{equation}
    where $c_{1-\alpha}^{boot}$ is the empirical $1-\alpha$ quantile of
    $(\hat\mu^{*(b)}-\hat\mu)^T\hat V_{BD}^{-1}
    (\hat\mu^{*(b)}-\hat\mu)$;
    \item report the union of all structural sets implied by
    $(B,D)\in\mathcal E_{\alpha}^{boot}$.
\end{enumerate}

If the object of interest is only $\gamma$, report the projection of that union
onto the $\gamma$ axis.  Under only \eqref{eq:trigmmset-region-consolidated},
that projection can be one-sided or unbounded.  A finite interval for $\gamma$
requires an extra restriction, for example a bounded range for the common-factor
loading $\beta=\alpha-\gamma$.

\subsection{Moment-inequality and test-inversion bootstrap}
\label{subsec:moment-ineq-bootstrap}

Many set-identified confidence procedures can be expressed as test inversion.
For each candidate value $w$ of $\theta$, define a null hypothesis
$H_0:\theta=w$.  Compute a statistic $T_n(w)$ that measures whether the sample
moments or inequalities are violated at $w$.  Retain $w$ if the test does not
reject.  The confidence set is
\begin{equation}
\mathcal C_{1-\alpha}=\{w:T_n(w)\le c_{1-\alpha}(w)\}.
\end{equation}

For a scalar moment-inequality representation, a common statistic is
\begin{equation}
T_n(w)=\max_{j\le J}
\frac{\sqrt n\,\bar m_j(w)}{\hat\sigma_j(w)},
\end{equation}
where the null requires $\E[m_j(w)]\le0$ for all $j$.  The critical value can be
computed by generalized moment selection, subsampling, or an $m$-out-of-$n$
bootstrap.  The R implementation is necessarily custom unless one can adapt a
specialized package from a different application.

The generic R loop is:

\begin{verbatim}
theta_grid <- seq(theta_min, theta_max, length.out = 1000)
keep <- logical(length(theta_grid))

for (k in seq_along(theta_grid)) {
  w <- theta_grid[k]

  # 1. compute sample moment inequalities mbar(w)
  # 2. compute a studentized test statistic Tn_w
  # 3. get a critical value by bootstrap, m-out-of-n bootstrap, or subsampling
  # 4. retain w if not rejected
  keep[k] <- (Tn_w <= c_alpha_w)
}

ci_theta <- range(theta_grid[keep])
\end{verbatim}

This approach is more work than bootstrapping endpoints, but it is conceptually
clean when the scalar coefficient is a projection of a higher-dimensional
identified set, as in \texttt{trigmmset}.

\subsection{Subsampling and $m$-out-of-$n$ bootstrap}
\label{subsec:subsampling-m-out-n}

Set endpoints are often nonsmooth functions of sample moments: a discriminant
can cross zero, a constraint can switch from slack to binding, or a projection
can move from bounded to unbounded.  In such cases, the ordinary $n$-out-of-$n$
bootstrap may be unreliable.  Two robust alternatives are:
\begin{enumerate}
    \item \emph{subsampling}: draw subsets of size $b$ without replacement,
    with $b\to\infty$ and $b/n\to0$;
    \item \emph{$m$-out-of-$n$ bootstrap}: draw $m<n$ observations with
    replacement, again with $m\to\infty$ and $m/n\to0$.
\end{enumerate}
For implementation, replace the line
\begin{verbatim}
ii <- sample.int(n, n, replace = TRUE)
\end{verbatim}
by either
\begin{verbatim}
ii <- sample.int(n, b, replace = FALSE)  # subsampling
\end{verbatim}
or
\begin{verbatim}
ii <- sample.int(n, m, replace = TRUE)   # m-out-of-n bootstrap
\end{verbatim}
The statistic must then be scaled appropriately.  For an endpoint statistic,
compare $\sqrt b(\hat L_b-\hat L)$ or $\sqrt m(\hat L_m^*-\hat L)$ to
$\sqrt n(\hat L-L_0)$.  In applied work, report sensitivity across a few values
of $b$ or $m$ rather than relying on one arbitrary tuning value.

\subsection{Recovering confidence intervals for the other coefficients}
\label{subsec:other-coefficients-inference}

Because $\beta_1(w)=\beta_1^R-\beta_2^Rw$, a set or interval for $\theta$ maps
into a set for each coefficient.  If $\theta\in[L,U]$, the $k$th coefficient is
in
\begin{equation}
\left\{\beta_{1,k}^R-\beta_{2,k}^R w:w\in[L,U]\right\}.
\end{equation}
If $\beta_{2,k}^R>0$, this interval is
\begin{equation}
[\beta_{1,k}^R-\beta_{2,k}^R U,
 \beta_{1,k}^R-\beta_{2,k}^R L].
\end{equation}
If $\beta_{2,k}^R<0$, the endpoints are reversed.  In a bootstrap, the safest
procedure is to recompute the residualization and recovery map in each draw:
\begin{equation}
\beta_{1,k}^{*(b)}(w)=
\beta_{1,k}^{R,*(b)}-\beta_{2,k}^{R,*(b)}w.
\end{equation}
This automatically accounts for the sampling uncertainty in the reduced-form
coefficients and in the identified interval endpoints.

\subsection{Recommended reporting template}
\label{subsec:reporting-template}

For the scalar triangular model, a clear empirical table should report:
\begin{enumerate}
    \item the OLS slope of $W_1$ on $W_2$;
    \item the exact Lewbel point estimate, if the exact covariance condition is
    maintained;
    \item first-stage relevance for the generated instrument, e.g.
    $\widehat{\Cov}(Z,W_2^2)$ and a heteroskedasticity test;
    \item the Lewbel bounded-invalidity interval for several values of $\tau$;
    \item bootstrap endpoint intervals for each reported set;
    \item if using the common-factor method, the \texttt{trigmm} point estimate
    or the \texttt{trigmmset}-style joint region and its projection for
    $\theta$;
    \item a statement of the coverage target: endpoint coverage, identified-set
    coverage, or true-parameter coverage.
\end{enumerate}
A short but rigorous sentence is: ``The bootstrap resamples observations and
re-estimates the residualizations, generated instruments, and set endpoints in
each replication; the reported interval is an outer endpoint interval for the
identified set.''

\section{Practical conclusions for software choice}
\label{sec:software-conclusions}

The most direct software choice depends on the identification target.
For Lewbel point identification in R, use \texttt{REndo::}\newline\texttt{hetErrorsIV()} or
manually construct the generated instrument and use \texttt{ivreg}.  For
bootstrap inference, manually wrap the whole construction in \texttt{boot}.  For
clustered point-identified IV inference, use \texttt{fwild}\newline\texttt{clusterboot} or
\texttt{clusterSEs} after constructing the generated instrument.

For Lewbel's Theorem 3 set-identified interval, no standard R package was found
that automates the method.  The necessary computation is short: compute
residuals, solve the quadratic for the two endpoints, and bootstrap the full
procedure.  The code in Section \ref{subsec:lewbel-theorem3-bootstrap} is the
most direct R implementation.

For the Lee--Lewbel--Schennach--Zhang common-factor model, Stata remains the
direct implementation: \texttt{trigmm} for point identification and
\texttt{trigmmset} for the second-moment set bound.  In R, \texttt{gmm} can be
used to code the nonlinear moments for point identification, while a
\texttt{trigmmset}-style set analysis requires custom code for the $(B,D)$
confidence region and the projection step.

For publication-quality partial-identification inference, the most credible
approach is to report both a transparent endpoint-bootstrap interval and a
sensitivity check using subsampling or $m$-out-of-$n$ bootstrap when endpoints
are unstable.  If the scalar parameter is a projection of a multidimensional
moment-inequality set, use test inversion or calibrated projection rather than a
naive percentile interval.

\begin{thebibliography}{99}

\bibitem[Baum and Lewbel(2019)]{baumlewbel2019}
Baum, Christopher F., and Arthur Lewbel. 2019. ``Advice on Using
Heteroskedasticity-Based Identification.'' \emph{Stata Journal} 19(4): 757--767.

\bibitem[Comon(1994)]{comon1994}
Comon, Pierre. 1994. ``Independent Component Analysis, a New Concept?''
\emph{Signal Processing} 36(3): 287--314.

\bibitem[Gourieroux, Monfort, and Renne(2017)]{gourieroux2017}
Gourieroux, Christian, Alain Monfort, and Jean-Paul Renne. 2017. ``Statistical
Inference for Independent Component Analysis: Application to Structural VAR
Models.'' \emph{Journal of Econometrics} 196(1): 111--126.

\bibitem[Guay(2021)]{guay2021}
Guay, Alain. 2021. ``Identification of Structural Vector Autoregressions Through
Higher Unconditional Moments.'' \emph{Journal of Econometrics} 225(1): 27--46.

\bibitem[Hyvarinen(1999)]{hyvarinen1999}
Hyvarinen, Aapo. 1999. ``Fast and Robust Fixed-Point Algorithms for Independent
Component Analysis.'' \emph{IEEE Transactions on Neural Networks} 10(3): 626--634.

\bibitem[Keweloh(2021)]{keweloh2021}
Keweloh, Sascha A. 2021. ``A Generalized Method of Moments Estimator for
Structural Vector Autoregressions Based on Higher Moments.'' \emph{Journal of
Business \& Economic Statistics} 39(3): 772--782.

\bibitem[Lanne, Meitz, and Saikkonen(2017)]{lanne2017}
Lanne, Markku, Mika Meitz, and Pentti Saikkonen. 2017. ``Identification and
Estimation of Non-Gaussian Structural Vector Autoregressions.'' \emph{Journal of
Econometrics} 196(2): 288--304.

\bibitem[Lee et al.(2026)]{leeetal2026}
Lee, Heejun, Arthur Lewbel, Susanne M. Schennach, and Linqi Zhang. 2026.
``Instrument-Free Estimation of Triangular Equation Systems with the
\texttt{trigmm} Command.'' \emph{Stata Journal} 26(1): 68--89.

\bibitem[Lewbel(2012)]{lewbel2012}
Lewbel, Arthur. 2012. ``Using Heteroscedasticity to Identify and Estimate
Mismeasured and Endogenous Regressor Models.'' \emph{Journal of Business \&
Economic Statistics} 30(1): 67--80.

\bibitem[Lewbel, Schennach, and Zhang(2024)]{lsz2024}
Lewbel, Arthur, Susanne M. Schennach, and Linqi Zhang. 2024. ``Identification of
a Triangular Two Equation System without Instruments.'' \emph{Journal of Business
\& Economic Statistics} 42(1): 14--25.

\bibitem[Lewis(2021)]{lewis2021}
Lewis, Daniel J. 2021. ``Identifying Shocks via Time-Varying Volatility.''
\emph{Review of Economic Studies} 88(6): 3086--3124.

\bibitem[Lewis(2022)]{lewis2022}
Lewis, Daniel J. 2022. ``Robust Inference in Models Identified via
Heteroskedasticity.'' \emph{Review of Economics and Statistics} 104(3): 510--524.

\bibitem[Lewis(2025)]{lewis2025}
Lewis, Daniel J. 2025. ``Identification Based on Higher Moments in
Macroeconometrics.'' \emph{Annual Review of Economics} 17: 665--693.

\bibitem[Rigobon(2003)]{rigobon2003}
Rigobon, Roberto. 2003. ``Identification Through Heteroskedasticity.''
\emph{Review of Economics and Statistics} 85(4): 777--792.

\bibitem[Sentana and Fiorentini(2001)]{sentana2001}
Sentana, Enrique, and Gabriele Fiorentini. 2001. ``Identification, Estimation and
Testing of Conditionally Heteroskedastic Factor Models.'' \emph{Journal of
Econometrics} 102(2): 143--164.



\bibitem[Andrews and Soares(2010)]{andrewssoares2010}
Andrews, Donald W. K., and Gustavo Soares. 2010. ``Inference for Parameters
Defined by Moment Inequalities Using Generalized Moment Selection.''
\emph{Econometrica} 78(1): 119--157.

\bibitem[Canty and Ripley(2024)]{bootpackage}
Canty, Angelo, and Brian Ripley. 2024. \texttt{boot}: Bootstrap Functions.
R package.  Based on Davison and Hinkley (1997).  URL:
\url{https://cran.r-project.org/package=boot}.

\bibitem[Chausse(2010)]{gmmpackage}
Chausse, Pierre. 2010. ``Computing Generalized Method of Moments and Generalized
Empirical Likelihood with R.'' \emph{Journal of Statistical Software} 34(11):
1--35.

\bibitem[Fox, Kleiber, and Zeileis(2024)]{ivregpackage}
Fox, John, Christian Kleiber, and Achim Zeileis. 2024. \texttt{ivreg}:
Instrumental-Variables Regression by 2SLS, 2SM, or 2SMM, with Diagnostics. R
package. URL: \url{https://cran.r-project.org/package=ivreg}.

\bibitem[Gui, Meierer, and Algesheimer(2023)]{rendo2023}
Gui, Raluca, Markus Meierer, and Rene Algesheimer. 2023. ``REndo: Internal
Instrumental Variables to Address Endogeneity.'' \emph{Journal of Statistical
Software} 107(3): 1--43.

\bibitem[Imbens and Manski(2004)]{imbensmanski2004}
Imbens, Guido W., and Charles F. Manski. 2004. ``Confidence Intervals for
Partially Identified Parameters.'' \emph{Econometrica} 72(6): 1845--1857.

\bibitem[Kaido, Molinari, and Stoye(2019)]{kaidomolinaristoye2019}
Kaido, Hiroaki, Francesca Molinari, and Jorg Stoye. 2019. ``Confidence Intervals
for Projections of Partially Identified Parameters.'' \emph{Econometrica} 87(4):
1397--1432.

\bibitem[MacKinnon, Nielsen, and Webb(2022)]{mackinnon2022}
MacKinnon, James G., Morten Nielsen, and Matthew D. Webb. 2022. ``Fast and
Reliable Jackknife and Bootstrap Methods for Cluster-Robust Inference.''
\emph{Journal of Applied Econometrics} 37(5): 944--967.

\bibitem[Romano and Shaikh(2010)]{romanosh2010}
Romano, Joseph P., and Azeem M. Shaikh. 2010. ``Inference for the Identified Set
in Partially Identified Econometric Models.'' \emph{Econometrica} 78(1):
169--211.

\bibitem[Stoye(2009)]{stoye2009}
Stoye, Jorg. 2009. ``More on Confidence Intervals for Partially Identified
Parameters.'' \emph{Econometrica} 77(4): 1299--1315.

\end{thebibliography}

\end{document}
